% =============================================================
%  INTRODUCTION
% =============================================================
\section{Introduction}

Intermolecular RNA--RNA interactions underpin a wide range of regulatory mechanisms, and predicting them at genome- or transcriptome-wide scale is a recurring task in RNA biology.
The problem is particularly demanding for synthetic small interfering RNAs (siRNAs), where unintended off-target binding can confound RNAi knockdown experiments and limit therapeutic application.

RIsearch2~\citep{alkan2017risearch2} is a tool designed for large-scale screens of near-complementary RNA--RNA interactions: it combines a suffix-array index of the target with a seed-and-extend dynamic-programming step that uses the simplified energy model introduced by RIsearch~\citep{wenzel2012risearch}.
The companion siRNA off-target pipeline distributed alongside RIsearch2 turns these predictions into per-transcript binding probabilities and a transcriptome-wide off-targeting potential through a partition-function model that integrates transcript abundance and binding-site accessibility.

Despite the algorithm's continued relevance, the software ecosystem has aged less well: RIsearch2 ships as a C source tarball with a custom Makefile; the off-target pipeline is written in Python~2.7, long past end of life; and neither component exposes a library API suitable for embedding in modern Python-native analyses.

We introduce new versions of both tools, redesigned for improved usability, interoperability with modern analysis frameworks, and easier installation.
\textbf{RIsearch}~(v3) is a Rust reimplementation of the RIsearch2 search algorithm---ranking among the fastest tools for genome-wide RNA--RNA interaction scans---with parity-tested behaviour, a redesigned command-line interface, and PyO3 Python bindings that expose search results as native Polars~\citep{polars} DataFrames.
\textbf{RIOT} (RIsearch siRNA Off-Target pipeline) is a rewritten siRNA off-target pipeline that consumes those bindings directly; it thus doubles as a reference example for users wishing to embed RIsearch in their own analyses.

% Figure 1 declared here so the two-column float lands at the top of page 2.
\begin{figure*}[t]
\centering
\resizebox{\linewidth}{!}{% RIsearch + RIOT overview — TikZ source (Figure 1).
% A two-lane install -> execute -> output flow, blocks aligned on one row per lane.
%   Lane 1: RIsearch  (Crates.io / PyPI install, SA seed-and-extend, generic output).
%   Lane 2: RIOT      (PyPI install, RIsearch-backed pipeline, final off-target file).
% Required libraries (already loaded in main.tex for summary.tex):
%   \usetikzlibrary{arrows.meta,positioning,backgrounds,fit,calc}

% ── Palette (matches the package's two-tool colour identity) ─────────────────
\definecolor{CoreBlue}{RGB}{179,208,232}     % RIsearch fill
\definecolor{CoreBorder}{RGB}{68,114,196}    % RIsearch border
\definecolor{PipeGreen}{RGB}{198,224,180}    % RIOT fill
\definecolor{PipeBorder}{RGB}{83,153,51}     % RIOT border
\definecolor{CodeBg}{RGB}{246,246,248}       % code-snippet background
\definecolor{CodeBd}{RGB}{200,200,205}       % code-snippet border
\definecolor{ChanBg}{RGB}{255,255,255}       % install-channel chip
\definecolor{InkGray}{RGB}{70,70,70}
\definecolor{VrnaGold}{RGB}{229,178,72}      % ViennaRNA accent
\definecolor{NclsPlum}{RGB}{150,110,170}     % NCLS accent

\begin{tikzpicture}[
  font=\sffamily,
  >={Stealth[length=2.6mm]},
  laneband/.style={
    rounded corners=3pt, fill=#1, text=white, anchor=south west,
    minimum height=6mm, inner xsep=6pt, font=\sffamily\bfseries\small,
  },
  stage/.style={
    rounded corners=6pt, draw=#1, line width=0.9pt, fill=#1!18,
    align=center, inner sep=5pt, font=\sffamily\small\bfseries, text=InkGray,
  },
  info/.style={
    rounded corners=3pt, draw=#1!55, line width=0.5pt, fill=#1!8,
    align=left, inner sep=5pt, font=\sffamily\scriptsize, text=black,
  },
  code/.style={
    rounded corners=3pt, draw=CodeBd, line width=0.5pt, fill=CodeBg,
    align=left, inner sep=4pt, font=\ttfamily\scriptsize, text=black,
  },
  chip/.style={
    rounded corners=3pt, draw=#1, line width=0.7pt, fill=ChanBg,
    align=center, inner sep=3pt, font=\sffamily\scriptsize\bfseries, text=#1,
  },
  flow/.style={->, line width=1pt, draw=#1},
  caplbl/.style={font=\sffamily\scriptsize, text=InkGray, align=center},
]

% column centres (shared by both lanes so the two rows line up vertically)
\def\colA{0}
\def\colB{6.2}
\def\colC{12.4}

% =====================================================================
%  LANE 1 — RIsearch  (all blocks on one row, tops aligned at y = 0)
% =====================================================================
% ---- (1a) Install: two channels stacked vertically (no overlap) -----
\node[stage=CoreBorder, text width=3.4cm, anchor=north] (rs-install) at (\colA,0)
  {Install (Rust core)};
\node[chip=CoreBorder, anchor=north, text width=1.8cm]
  (rs-crates) at ([yshift=-4pt]rs-install.south) {Crates.io};
\node[code, anchor=north, text width=3.35cm]
  (rs-crates-cmd) at ([yshift=-2pt]rs-crates.south)
  {\$ cargo install risearch\\ \$ risearch --help};
\node[chip=CoreBorder, anchor=north, text width=1.8cm]
  (rs-pypi) at ([yshift=-5pt]rs-crates-cmd.south) {PyPI};
\node[code, anchor=north, text width=3.35cm]
  (rs-pypi-cmd) at ([yshift=-2pt]rs-pypi.south)
  {\$ pip install risearch\\ >>> import risearch};
\node[draw=CoreBorder!45, line width=0.6pt, rounded corners=6pt,
      fit=(rs-install)(rs-crates-cmd)(rs-pypi-cmd), inner sep=4pt] (rs-install-box) {};

% ---- (1b) Execute: conceptual description (no commands) -------------
\node[stage=CoreBorder, text width=4.55cm, anchor=north] (rs-run) at (\colB,0)
  {Index \& search};
\node[info=CoreBorder, anchor=north, text width=4.55cm]
  (rs-run-info) at ([yshift=-4pt]rs-run.south)
  {\textbullet~Suffix-array (SA) index of targets\\[2pt]
   \textbullet~Seed-and-extend DP extension\\[2pt]
   \textbullet~Simplified RIsearch energy model};
\node[draw=CoreBorder!45, line width=0.6pt, rounded corners=6pt,
      fit=(rs-run)(rs-run-info), inner sep=4pt] (rs-run-box) {};

% ---- (1c) Output: generic, illustrative (not the real schema) -------
\node[stage=CoreBorder, text width=5.0cm, anchor=north] (rs-out) at (\colC,0)
  {RIsearch output};
\node[code, anchor=north, text width=5.0cm]
  (rs-out-tbl) at ([yshift=-4pt]rs-out.south)
  {query   target   start   end   energy\\[1pt]
   q1      chrX      881    901   -28.4\\
   q1      chrY      233    253   -24.1\\[3pt]
   \# TSV output};
\node[draw=CoreBorder!45, line width=0.6pt, rounded corners=6pt,
      fit=(rs-out)(rs-out-tbl), inner sep=4pt] (rs-out-box) {};

% horizontal lane arrows (both pinned to the run-header centre line)
\draw[flow=CoreBorder] (rs-install-box.east|-rs-run) -- (rs-run-box.west|-rs-run);
\draw[flow=CoreBorder] (rs-run-box.east|-rs-run) -- (rs-out-box.west|-rs-run);

% lane-1 label tab
\node[laneband=CoreBorder]
  (rs-tab) at ([yshift=4pt]rs-install-box.north west)
  {RIsearch \textemdash\ RNA--RNA interaction search core};

% bounding helper for clearance below lane 1
\node[fit=(rs-install-box)(rs-run-box)(rs-out-box), inner sep=0] (lane1) {};

% =====================================================================
%  LANE 2 — RIOT  (all blocks on one row, tops aligned below lane 1)
% =====================================================================
\coordinate (l2) at ([yshift=-2.0cm]lane1.south -| rs-install);

% ---- (2a) Install ---------------------------------------------------
\node[stage=PipeBorder, text width=4.1cm, anchor=north] (ro-install) at (l2)
  {Install (RIOT pipeline)};
\node[chip=PipeBorder, anchor=north, text width=1.8cm]
  (ro-pypi) at ([yshift=-4pt]ro-install.south) {PyPI};
\node[code, anchor=north, text width=4.1cm]
  (ro-install-cmd) at ([yshift=-2pt]ro-pypi.south)
  {\$ pip install riot\\[1pt]
   \# bundles: risearch wheel,\\
   \# ViennaRNA API, NCLS\\
   \# uv-locked, no binaries};
\node[draw=PipeBorder!45, line width=0.6pt, rounded corners=6pt,
      fit=(ro-install)(ro-install-cmd), inner sep=4pt] (ro-install-box) {};

% ---- (2b) Pipeline (built on RIsearch) ------------------------------
\node[stage=PipeBorder, text width=4.55cm, anchor=north]
  (ro-run) at ([xshift=\colB cm]ro-install.north) {Run off-target pipeline};
\node[code, anchor=north, text width=4.55cm]
  (ro-run-cmd) at ([yshift=-4pt]ro-run.south)
  {\$ riot off-targets \textbackslash\\
   \ \ \ -s sirna.fa --genome hg38.fa \textbackslash\\
   \ \ \ -t annot.gtf -a access/ \textbackslash\\
   \ \ \ --temperature 37 -o off.tsv\\[2pt]
   \# or: riot -c run.yaml};
\node[stage=CoreBorder, text width=4.4cm, anchor=north, font=\sffamily\scriptsize\bfseries]
  (st-search) at ([yshift=-4pt]ro-run-cmd.south) {1. RIsearch index + search (in-process)};
\node[stage=NclsPlum, text width=4.4cm, anchor=north, font=\sffamily\scriptsize\bfseries, fill=NclsPlum!18, draw=NclsPlum]
  (st-isect) at ([yshift=-3pt]st-search.south) {2. Intersect annotation \textbullet\ NCLS};
\node[stage=VrnaGold, text width=4.4cm, anchor=north, font=\sffamily\scriptsize\bfseries, fill=VrnaGold!22, draw=VrnaGold]
  (st-acc) at ([yshift=-3pt]st-isect.south) {3. Accessibility \textbullet\ ViennaRNA API};
\node[stage=PipeBorder, text width=4.4cm, anchor=north, font=\sffamily\scriptsize\bfseries]
  (st-prob) at ([yshift=-3pt]st-acc.south) {4. Boltzmann partition function $Z_s$};
\node[draw=PipeBorder!45, line width=0.6pt, rounded corners=6pt,
      fit=(ro-run)(ro-run-cmd)(st-search)(st-prob), inner sep=4pt] (ro-run-box) {};

\draw[flow=InkGray, line width=0.7pt] (st-search.south) -- (st-isect.north);
\draw[flow=InkGray, line width=0.7pt] (st-isect.south) -- (st-acc.north);
\draw[flow=InkGray, line width=0.7pt] (st-acc.south) -- (st-prob.north);

% ---- (2c) Final output file -----------------------------------------
\node[stage=PipeBorder, text width=5.0cm, anchor=north]
  (ro-out) at ([xshift=\colC cm]ro-install.north) {RIOT output};
\node[code, anchor=north, text width=5.0cm]
  (ro-out-tbl) at ([yshift=-4pt]ro-out.south)
  {sirna\_id   transcript\_id   P\_off\\[1pt]
   si\_001     ENST...0412      0.973\\
   si\_001     ENST...1097      0.118\\[3pt]
   \# TSV output};
\node[draw=PipeBorder!45, line width=0.6pt, rounded corners=6pt,
      fit=(ro-out)(ro-out-tbl), inner sep=4pt] (ro-out-box) {};

% horizontal lane arrows (both pinned to the run-header centre line)
\draw[flow=PipeBorder] (ro-install-box.east|-ro-run) -- (ro-run-box.west|-ro-run);
\draw[flow=PipeBorder] (ro-run-box.east|-ro-run) -- (ro-out-box.west|-ro-run);

% lane-2 label tab
\node[laneband=PipeBorder]
  (ro-tab) at ([yshift=4pt]ro-install-box.north west)
  {RIOT \textemdash\ RIsearch siRNA Off-Target pipeline};

% ---- "built on RIsearch" connector: straight segments, routed down the
%      clear corridor between the pipeline box and the off-target table so
%      it passes beside/under those boxes rather than through them.
\coordinate (cdrop) at ([yshift=-0.5cm]rs-out-box.south);            % just below Output
\coordinate (ccorr) at (9.2, 0 |- cdrop);                            % left into the corridor
\coordinate (cdown) at (9.2, 0 |- st-search.east);                   % straight down the corridor
\draw[flow=CoreBorder, line width=1pt, dashed]
  (rs-out-box.south) -- (cdrop) -- (ccorr) -- (cdown) -- (st-search.east);
\node[caplbl, fill=white, inner sep=1.5pt]
  at ([yshift=-0.3cm]$(cdrop)!0.5!(ccorr)$)
  {same engine, same process};

\end{tikzpicture}
}
\caption{%
  \textbf{Overview of RIsearch and RIOT.}
  \textbf{(Top, blue)} \textbf{RIsearch} is a Rust reimplementation of the RIsearch2 search algorithm, installable as a single static binary via \texttt{cargo install risearch} (Crates.io) or as a Python wheel via \texttt{pip install risearch} (PyPI).
  The core builds a suffix-array index of target sequences and runs a seed-and-extend dynamic-programming search using the simplified RIsearch energy model, producing interaction hits as a TSV file.
  \textbf{(Bottom, green)} \textbf{RIOT} (RIsearch siRNA Off-Target pipeline) is a modern Python~3 rewrite installable with a single \texttt{pip install riot}, which bundles the RIsearch wheel, the ViennaRNA~2.7.2 Python API, and an in-memory NCLS interval index with no external binaries required.
  RIOT orchestrates four stages: (\textit{i})~RIsearch index and search invoked natively via PyO3 bindings; (\textit{ii})~intersection of predicted binding sites with a transcriptome annotation using an NCLS index; (\textit{iii})~per-site RNA accessibility computed through the ViennaRNA Python API; and (\textit{iv})~a per-siRNA Boltzmann partition function yielding off-target probabilities, written to a TSV output file.
  The dashed arrow indicates that RIOT calls RIsearch within the same process via the Python bindings, with no subprocess or intermediate files.%
}
\label{fig:summary}
\end{figure*}

% =============================================================
%  RATIONALE FOR A RUST REIMPLEMENTATION
% =============================================================
\section{Rationale for a Rust Reimplementation}\label{sec:rationale}

What limits RIsearch2 today is not its algorithm but how that algorithm is delivered.
A reimplementation must meet three requirements at once: retain genome-scale throughput, so whole-transcriptome screens stay practical; install as a prebuilt dependency; and run as a safe Python library, so off-target search composes with surrounding analysis.
Each is routine alone; no single established language meets all three.
C reaches that throughput but has no standard way to deliver a prebuilt, installable dependency and cannot guarantee memory safety at a library boundary; Python distributes and embeds easily but in pure Python cannot reach that throughput.
Rust meets all three: it compiles to native code, installs as a prebuilt package, and exposes the algorithm through a memory-safe Python library.
That combination is why we built RIsearch in Rust.

% =============================================================
%  RISEARCH CORE IMPLEMENTATION
% =============================================================
\section{RIsearch Core Implementation}\label{sec:core}\label{subsec:pybindings}

RIsearch reimplements the RIsearch2 seed-and-extend pipeline.
Targets are concatenated and indexed in a suffix array, built in linear time by libsais~\citep{libsais}, and each query is scanned against the index for near-complementary seeds, permitting G--U wobble pairs and a configurable number of mismatches.
Seeds are extended in both directions with Gotoh's three-state affine-gap recurrence~\citep{gotoh1982improved}, scoring stacked pairs from a nearest-neighbour dinucleotide model and penalising bulges and internal loops.
The algorithm and energy model are unchanged, so predicted hybridisation energies match those of RIsearch2.

\subsection*{Easier to install}

RIsearch2 is distributed only as a C source tarball with a custom Makefile, so every user compiles it locally, which presumes a C toolchain and leaves build flags and dependency versions to each machine.
RIsearch, by contrast, installs as a prebuilt package from standard registries, either a precompiled PyPI wheel via \texttt{pip install risearch} or a \texttt{cargo install risearch} build from Crates.io with dependencies version-locked.
The result is a single self-contained binary that installs in one command and behaves identically across machines.

\subsection*{Easier to use}

Where RIsearch2 can return many overlapping alignments for one duplex, RIsearch reports a single lowest-energy alignment, collapsing the seeds that extend to the same site, and the reported hit set does not depend on the number of threads.
It validates its inputs up front and, on malformed data, reports a specific error that names the offending record, character, or row rather than failing obscurely.
The command line is rebuilt around task-oriented subcommands, and legacy RIsearch2 flags are still accepted so existing scripts run unchanged.

\subsection*{Easier to integrate}

RIsearch2 exposed no library API, so embedding it in a Python analysis meant spawning a subprocess and parsing its text output.
RIsearch instead ships PyO3 bindings in the same wheel: \texttt{risearch.index()} builds an on-disk index, \texttt{risearch.TargetRegistry.open()} loads it as a reusable memory-mapped handle, and \texttt{risearch.search()} returns hits as a Polars DataFrame through the Arrow C stream interface, with no intermediate files and no text parsing:

\begin{verbatim}
import risearch
risearch.index("targets.fa", "index/")
store = risearch.TargetRegistry.open("index/")
df = risearch.search(queries, store, matrix="t04")  # Polars DataFrame
\end{verbatim}

Searches release Python's global interpreter lock, so they run concurrently with other work in the same process, and the energy parameter set is selectable per call, defaulting to \texttt{t04}.

\subsection*{Easier to extend}

Energy models are data rather than code: each parameter set is declared in a manifest with its published provenance and validated at build time, so supporting new thermodynamic parameters is a data file rather than an engine change.
The engine itself is modular and open source, and is covered by an automated test suite that includes mutation testing, so contributors can extend it and re-verify its behaviour with confidence.

\subsection*{Performance}

On equivalent inputs, RIsearch is $2$--$3\times$ faster than RIsearch2 (Supplementary Figures~S1--S5; dataset, protocol, and full results in Supplementary Methods).
Its suffix array is also built in parallel across cores, about $2\times$ faster than RIsearch2's single-threaded libdivsufsort~\citep{libdivsufsort} construction on the hg38 genome (Supplementary Figure~S6).
The on-disk index is memory-mapped and read without deserialisation, so a multi-gigabyte genome index is available immediately and shared across processes by the operating system, rather than read into memory on every run.

% =============================================================
%  SIRNA OFF-TARGET DISCOVERY PIPELINE
% =============================================================

\section{RIOT: a Modernised siRNA Off-Target pipeline}\label{sec:pipeline}

Alongside the core, we distribute \textbf{RIOT} (RIsearch siRNA Off-Target pipeline), a complete reimplementation of the siRNA off-target pipeline of \citet{alkan2017risearch2} (Figure~\ref{fig:summary}).
Beyond porting the workflow from Python~2.7 to Python~3, the rewrite targets three practical barriers that limited adoption of the original tool: installation, day-to-day use, and integration into existing analyses.
The underlying four-stage thermodynamic model---search, annotation intersection, accessibility, and a per-siRNA Boltzmann partition function---is preserved unchanged and described in full in the Supplementary Methods.

\subsection*{Easier to install}

The original pipeline required a Python~2.7 interpreter, a separately compiled RIsearch2 C binary, and externally installed \texttt{RNAplfold} (ViennaRNA) and \texttt{bedtools}~\citep{quinlan2010bedtools} executables on the user's \texttt{PATH}.
RIOT removes every external binary: a single \texttt{pip install riot} brings in the RIsearch core as a pre-compiled wheel, the ViennaRNA~2.7.2 Python API~\citep{lorenz2011vienna} in place of the \texttt{RNAplfold} executable, and an in-memory NCLS interval index~\citep{ncls} in place of \texttt{bedtools}.
All transitive dependencies are resolved and pinned through a universal \texttt{uv} lockfile, so installations are reproducible and free of version conflicts across machines.

\subsection*{Easier to use}

The legacy pipeline was a single \texttt{argparse} script driven by a fixed set of command-line flags.
RIOT retains backward-compatible flags but adds a YAML configuration mode: every option---inputs, the RIsearch energy and seed thresholds, the saturation parameters $\alpha$, $\gamma$, and $\theta$, and the output format---can be declared in one file, validated at startup against a typed OmegaConf schema.
A single configuration file captures an entire run, making analyses straightforward to reproduce and to share between collaborators.
For large screens, RIOT integrates natively with Slurm~\citep{yoo2003slurm}: one orchestrator submits the indexing, accessibility, and off-target stages as dependent batch jobs with per-stage resource requests, replacing hand-written submission scripts.

A single \texttt{-{}-temperature} option (default $37\,^\circ$C) propagates to both the ViennaRNA accessibility calculation and the Boltzmann partition function, so that hybridisation and opening energies are always evaluated on a consistent thermodynamic scale---directly replicating the intent of the original RIsearch model~\citep{wenzel2012risearch} while exposing the folding condition as a first-class parameter (Supplementary Methods).

\subsection*{Easier to integrate}

Because the legacy pipeline communicated through files and subprocesses, embedding it in a larger analysis meant shelling out to a script and parsing its text output.
RIOT instead exposes its stages as an importable Python API: the search, intersection, accessibility, and probability services can each be called directly and return native Polars DataFrames, so a user can drop any stage into an existing Python workflow and pass results to other Python-native tools in a single process, with no intermediate files and no text parsing.
Built on the RIsearch core bindings (Section~\ref{subsec:pybindings}), RIOT thus serves both as a turnkey command-line tool and as a library component of bespoke off-target analyses.

\subsection*{Performance}

Benchmarked on an AMD EPYC 9554 (64 cores, 32 workers) using subsets of the genome-wide siRNA efficacy library of \citet{huesken2005design} searched against the GRCh38/hg38 reference genome~\citep{schneider2017grch38} and intersected against an H1299 RNA-seq expression profile (dataset, genome indexing, expression annotation, and benchmark protocol detailed in Supplementary Methods), RIOT completes the intersection, accessibility annotation, probability computation, and output stages $7.7{-}10.7\times$ faster than the legacy Python~2.7 implementation across batch sizes of 500--2{,}000 siRNAs (five randomised trials per batch size; Supplementary Figure~S7).
Peak memory plateaus at approximately 7\,GB regardless of batch size, whereas the legacy pipeline grows from $\sim$8.5\,GB to $\sim$10\,GB as batch size increases (Supplementary Methods).

Separately, the one-time genome-wide accessibility pre-computation---folding the entire hg38 reference at $W{=}80$, $L{=}40$, $u{=}30$ with 64 workers---runs $\sim$34\% faster than the legacy \texttt{RNAplfold} script (8\,h\,46\,m vs.\ 13\,h\,12\,m) and writes profiles $\sim$42\% smaller on disk (111.8\,GB vs.\ 192.6\,GB), at the cost of higher but still modest peak memory ($\sim$3.6\,GB vs.\ $\sim$309\,MB; Supplementary Figure~S8 and Supplementary Table~S1).

% =============================================================
%  FINAL REMARKS
% =============================================================
\section{Final Remarks}

% 1 short paragraph: what's available, who benefits, what's next (pooled-siRNA
% off-targeting potential, additional energy models, further integration targets).
